\scanpage{101}
\mspage{3.6}
% archive sheet number: 100
Corollary 3.9: Let $R\handwritten{\subset}S$ be a local algebra of the integrally closed local rings such that the quotients field extension $k\handwritten{\subset}K$ is galois. Then $r(S/R)={}$ the number of minimal primes $\underline p$ in $S^f$ such that $\underline p\cap R=0={}$ the order of the inertia group of $S$ over $R$.

\uline{Proof}: Let $U$ be the inertial ring of $S$ over $R$. Since $U$ is an unramified local $R$-algebra, we have by lemma 2.1 that $S_R^f\handwritten{\approx}S_U^f$. However $U$ is a transporter for $S$ over $R$ and hence $S_R^f\handwritten{\approx}S_U^f=S\mathbin{\ooalign{\hfil$\mathrm{x}$\hfil\cr\hfil$\handwritten{\bigcirc}$\hfil\cr}}_U S$. For a prime ideal $\underline p$ in $S_R^f$ it is clear that $\underline p\cap R=0$ if and only if $\underline p\cap U=0$, since $U$ is a localization of an integral extension of $R$. Therefore the minimal prime ideals $\underline p$ in $S_R^f$ with $\underline p\cap R=0$ are in one-one correspondence with the prime ideals in $(S\mathbin{\ooalign{\hfil$\mathrm{x}$\hfil\cr\hfil$\handwritten{\bigcirc}$\hfil\cr}}_U S)\mathbin{\ooalign{\hfil$\mathrm{x}$\hfil\cr\hfil$\handwritten{\bigcirc}$\hfil\cr}}_U K_I=K\mathbin{\ooalign{\hfil$\mathrm{x}$\hfil\cr\hfil$\handwritten{\bigcirc}$\hfil\cr}}_{K_I}K$ where $K_I$ is the field of quotients of $U$, i.e. the inertial field of $S$ over $R$. Now $K$ is a galois extension of $K_I$ with $G_I$ the inertia group of $S$ over $R$ as its galois group. We recall that there is one-one correspondence between the prime ideals in $K\mathbin{\ooalign{\hfil$\mathrm{x}$\hfil\cr\hfil$\handwritten{\bigcirc}$\hfil\cr}}_{K_I}K$ and the elements in $G_I$ under $\underline P(\handwritten{\sigma})=\mathrm{Ker}\,(\handwritten{\widetilde\sigma}=K\mathbin{\ooalign{\hfil$\mathrm{x}$\hfil\cr\hfil$\handwritten{\bigcirc}$\hfil\cr}}_{K_I}K\longrightarrow K)$ where $\handwritten{\widetilde\sigma}(x\mathbin{\ooalign{\hfil$\mathrm{x}$\hfil\cr\hfil$\handwritten{\bigcirc}$\hfil\cr}} y)=x\handwritten{\sigma}(y)$. Therefore the prime ideals $\underline p$ in $S_R^f\handwritten{\approx}S\mathbin{\ooalign{\hfil$\mathrm{x}$\hfil\cr\hfil$\handwritten{\bigcirc}$\hfil\cr}}_U S$ with $\underline p\cap R=0$ are in one-one correspondence with the elements in $G_I$ under $\underline p(\handwritten{\sigma})=\underline P(\handwritten{\sigma})\cap S_R^f=\underline P(\handwritten{\sigma})\cap(S\mathbin{\ooalign{\hfil$\mathrm{x}$\hfil\cr\hfil$\handwritten{\bigcirc}$\hfil\cr}}_U S)$. Since $G_I$ is the inertia group of $S$ over $R$, we have in particular that $\handwritten{\sigma}(S)=S$ for all $\handwritten{\sigma} \varepsilon G_I$, and hence $\underline p(\handwritten{\sigma})=\underline P(\handwritten{\sigma})\cap(S\mathbin{\ooalign{\hfil$\mathrm{x}$\hfil\cr\hfil$\handwritten{\bigcirc}$\hfil\cr}}_U S)=\mathrm{Ker}\,(\handwritten{\widetilde\sigma}:S\mathbin{\ooalign{\hfil$\mathrm{x}$\hfil\cr\hfil$\handwritten{\bigcirc}$\hfil\cr}}_U S\longrightarrow S)$ where $\handwritten{\widetilde\sigma}(x\mathbin{\ooalign{\hfil$\mathrm{x}$\hfil\cr\hfil$\handwritten{\bigcirc}$\hfil\cr}} y)=x\handwritten{\sigma}(y)$. Therefore $S_R^f/\underline p(\handwritten{\sigma})=S\mathbin{\ooalign{\hfil$\mathrm{x}$\hfil\cr\hfil$\handwritten{\bigcirc}$\hfil\cr}}_U S/\underline p(\handwritten{\sigma})=S$ and hence $e(S_R^f/\underline p(\handwritten{\sigma})\mid S)=1$ for all $\handwritten{\sigma} \varepsilon G_I$. Therefore $r(S/R)=\handwritten{\sum}_{\varepsilon G_I}e(S^f/\underline p(\handwritten{\sigma})\mid S)={}$ the number of prime ideals $\underline p$ in $S_R^f$ with $\underline p\cap R=0={}$ the order of $G_I$.

Proposition 3.10: Let $R$ be an integrally closed local domain with quotient field $k$. Let $K$ be a finite algebraic extension of $k$ and $S$ the
